\section{Trabajos relacionados}
U-Net se ha consolidado como la arquitectura base de segmentación biomédica por su simetría encoder-decoder y sus conexiones de salto, que permiten recuperar detalle espacial perdido durante la compresión \citep{ronneberger2015unet}. En segmentación retiniana, esta propiedad es particularmente importante porque los vasos finos ocupan pocos píxeles y desaparecen con facilidad cuando la representación se vuelve demasiado gruesa.

Attention U-Net extiende esta idea introduciendo compuertas que reponderan las \emph{skip connections} con información contextual del decoder \citep{oktay2018attention}. Esta modificación resulta útil cuando el fondo presenta regiones oscuras, reflejos o ruido que compiten con los vasos de baja intensidad. Por otro lado, U-Net++ incrementa la capacidad de agregación multiescala mediante conexiones densas anidadas entre niveles del decoder \citep{zhou2019unetpp}, reduciendo el desajuste semántico entre mapas superficiales y profundos.

En el dominio específico de vasos retinianos, la literatura reporta que el principal cuello de botella sigue siendo el desbalance extremo entre píxeles de fondo y de vaso, además del cambio de dominio entre centros clínicos \citep{fraz2012review,staal2004drive,hoover2000stare}. Por ello, además de la arquitectura, las decisiones de preprocesamiento y de función de pérdida tienen un impacto directo en el rendimiento. En este proyecto se comparan BCE, Dice y pérdidas combinadas, junto con una estrategia de adaptación basada en CLAHE \citep{pizer1987clahe} y TTA.
